\appendix

\section{Getting the source code}
The code for this thesis is hosted on github at
\url{https://github.com/svbrodersen/Speciale}. This repository contains the
work done for the thesis including the report itself, and the different
submodule projects used throughout the thesis. Further, to reproduce the
results a podman image is provided to get a container capable of building and
running the demo examples. The container is hosted at
\url{ghcr.io/svbrodersen/speciale:latest}. To help with the setup, the root of
the repository contains a README.md file with instructions on how to build and
run the demos along with how to setup the container. To clone the repository
run the following commands:

\begin{lstlisting}[language=bash]
git clone https://github.com/svbrodersen/Speciale.git
git submodule update --init --recursive
\end{lstlisting}
As the project has quite some large submodules, beware that this might take up
a significant amount of space, and could take some time.

The primary source code for the modified FreeRTOS kernel is hosted at:
\url{https://github.com/svbrodersen/FreeRTOS-Kernel/tree/main}, and it should be
pulled automatically following the submodule recursive update.
